% Squelette du working paper. Chaque chiffre, tableau et figure provient d'un script de scripts/.
% Compilation : make paper (latexmk -pdf). Les tableaux sont inclus depuis tables/*.tex, les figures
% depuis figures/*.pdf ; rien n'est saisi à la main.
\documentclass[11pt,a4paper]{article}
\usepackage[utf8]{inputenc}
\usepackage[T1]{fontenc}
\usepackage[french,english]{babel}
\usepackage{lmodern,microtype}
\usepackage[margin=2.5cm]{geometry}
\usepackage{amsmath,amssymb,booktabs,graphicx,setspace,threeparttable,longtable}
\usepackage[hidelinks]{hyperref}
\usepackage[authoryear,round]{natbib}
\graphicspath{{../figures/}}
\onehalfspacing

\title{Smartphones, réseaux sociaux, applications de rencontre et fécondité :\\
y a-t-il un effet causal au-delà des adolescentes ?}
\author{Working paper}
\date{\today\ --- version de travail, ne pas citer}

\begin{document}
\selectlanguage{french}
\maketitle

\begin{abstract}
\selectlanguage{english}
\noindent\textbf{Abstract.} [À rédiger à l'Étape 4, à partir des seuls résultats du papier.]
\selectlanguage{french}

\medskip\noindent\textbf{Résumé.} [À rédiger à l'Étape 4.]
\end{abstract}

\section{Introduction}
% question, contribution, résultat principal en une phrase

\section{Positionnement}
% ≤ 600 mots ; références vérifiées uniquement (references.bib)

\section{Cadre conceptuel et hypothèses}
% reprend docs/preregistration.md (hash du commit de gel : à insérer)

\section{Données}
% sources, construction, descriptives, limites de mesure ; tables/t_sample.tex ; docs/data_log.md

\section{Stratégie empirique}

\section{Résultats : effet causal par âge}
% figures/f_event_*.pdf ; tables/t_main_*.tex — aucun résultat externe dans les sections 6 à 11

\section{Résultats : décomposition du canal}

\section{Robustesse}

\section{Discussion et mécanismes}

\section{Limites}

\section{Conclusion et critères de révision}
% (i) statut du smartphone comme facteur (premier ordre / second ordre / indéterminé) selon la règle préenregistrée ;
% (ii) seuil de révision à la hausse ; (iii) ce qui ferait rétrograder l'hypothèse.

\bibliographystyle{plainnat}
\bibliography{references}

\appendix
\section{Tableaux complets}
\section{Dictionnaire des variables}
\section{Sources exactes}
% généré depuis docs/data_log.md par scripts/06_tables.py
\section{Pays examinés et exclus}
% préregistration §4.3

\end{document}
